% =====================================================================
%  MODÈLE de rapport de sécurité - à copier dans <classe>/<slug>/<slug>.tex
%
%  Compilation depuis le dossier du rapport :
%      latexmk -pdf <slug>.tex        (recommandé)
%      # ou : pdflatex <slug>.tex  (deux fois, pour la table des matières)
%
%  Le chemin \input suppose l'arborescence <classe>/<slug>/ (deux niveaux
%  sous la racine du dépôt). Les captures vont dans ./images/.
% =====================================================================
\documentclass[11pt,a4paper]{article}
% =====================================================================
%  Préambule commun - Rapports de sécurité
%  Inclus par chaque rapport après \documentclass, avant \begin{document} :
%      \documentclass[11pt,a4paper]{article}
%      % =====================================================================
%  Préambule commun - Rapports de sécurité
%  Inclus par chaque rapport après \documentclass, avant \begin{document} :
%      \documentclass[11pt,a4paper]{article}
%      % =====================================================================
%  Préambule commun - Rapports de sécurité
%  Inclus par chaque rapport après \documentclass, avant \begin{document} :
%      \documentclass[11pt,a4paper]{article}
%      \input{../../templates/preambule.tex}
%      \begin{document} ... \end{document}
%
%  Moteur : pdflatex (compiler deux fois pour la table des matières, ou
%  utiliser latexmk -pdf).
% =====================================================================

\usepackage[T1]{fontenc}
\usepackage[utf8]{inputenc}
\usepackage{lmodern}
\usepackage[french]{babel}
\usepackage{csquotes}
\frenchsetup{StandardItemLabels=true}
\usepackage{amsmath}

\usepackage[margin=2.5cm]{geometry}
\usepackage{graphicx}
\usepackage{float}
\usepackage{xcolor}
\usepackage{array}
\usepackage{booktabs}
\usepackage{enumitem}
\usepackage{listings}
\usepackage{tcolorbox}
\tcbuselibrary{skins,breakable}
\usepackage[hidelinks]{hyperref}

% ---------------------------------------------------------------------
%  Palette
% ---------------------------------------------------------------------
\definecolor{codebg}{RGB}{245,246,248}
\definecolor{codeframe}{RGB}{210,214,220}
\definecolor{codekw}{RGB}{170,40,60}
\definecolor{codestr}{RGB}{30,110,70}
\definecolor{codecmt}{RGB}{120,125,135}
\definecolor{accent}{RGB}{30,70,120}
\definecolor{warnbg}{RGB}{255,247,230}
\definecolor{warnframe}{RGB}{224,168,64}
\definecolor{keybg}{RGB}{235,240,248}
\definecolor{keyframe}{RGB}{120,150,190}

% ---------------------------------------------------------------------
%  Listings (blocs de code encadrés, monospace)
% ---------------------------------------------------------------------
\lstdefinestyle{report}{
  backgroundcolor=\color{codebg},
  frame=single,
  rulecolor=\color{codeframe},
  basicstyle=\ttfamily\small,
  keywordstyle=\color{codekw}\bfseries,
  stringstyle=\color{codestr},
  commentstyle=\color{codecmt}\itshape,
  numbers=none,
  breaklines=true,
  breakatwhitespace=false,
  showstringspaces=false,
  columns=fullflexible,
  keepspaces=true,
  xleftmargin=8pt,
  xrightmargin=8pt,
  aboveskip=10pt,
  belowskip=10pt,
  literate=%
    {é}{{\'e}}1 {è}{{\`e}}1 {ê}{{\^e}}1 {à}{{\`a}}1 {ç}{{\c c}}1
    {É}{{\'E}}1 {î}{{\^i}}1 {ï}{{\"i}}1 {ô}{{\^o}}1 {û}{{\^u}}1
    {ù}{{\`u}}1 {œ}{{\oe}}1 {»}{{\guillemotright}}1 {«}{{\guillemotleft}}1
}
\lstset{style=report}

% ---------------------------------------------------------------------
%  Captures d'écran : image + légende en italique (sans « Figure N »)
%  \shot{fichier.png}{légende}
%  Si l'image manque, un cadre de remplacement est affiché.
% ---------------------------------------------------------------------
\newcommand{\shot}[2]{%
  \begin{center}%
    \IfFileExists{images/#1}%
      {\includegraphics[width=0.92\linewidth]{images/#1}}%
      {\setlength{\fboxsep}{0pt}\fcolorbox{codeframe}{codebg}{%
         \parbox[c][3.2cm][c]{0.92\linewidth}{\centering\ttfamily\small
         images/#1\\[4pt]\normalfont\itshape(capture à ajouter)}}}%
    \\[5pt]{\itshape\small #2}%
  \end{center}%
}

% ---------------------------------------------------------------------
%  Encadrés « points clés » (bleu) et « attention » (orange)
% ---------------------------------------------------------------------
\newtcolorbox{keybox}[1][]{
  breakable, colback=keybg, colframe=keyframe, boxrule=0.6pt,
  arc=2pt, left=8pt, right=8pt, top=6pt, bottom=6pt, #1}
\newtcolorbox{warnbox}[1][]{
  breakable, colback=warnbg, colframe=warnframe, boxrule=0.6pt,
  arc=2pt, left=8pt, right=8pt, top=6pt, bottom=6pt, #1}

\setlist[itemize]{leftmargin=1.4em, itemsep=2pt, topsep=4pt}

% ---------------------------------------------------------------------
%  Page de garde + table des matières
%  \makecover{titre}{sous-titre}{lignes méta}{auteur}{date}
%  Les lignes méta s'écrivent avec \metarow{Étiquette}{Valeur}.
% ---------------------------------------------------------------------
\newcommand{\metarow}[2]{\textbf{#1} & #2 \\}
\newcommand{\makecover}[5]{%
  \begin{titlepage}
    \centering
    \vspace*{1.2cm}
    {\large\scshape Rapport de sécurité}\par
    \vspace{1.6cm}
    {\huge\bfseries #1\par}
    \vspace{0.5cm}
    {\Large #2\par}
    \vspace{2.2cm}
    \begin{tcolorbox}[width=0.86\textwidth, colback=codebg,
        colframe=codeframe, boxrule=0.8pt, arc=3pt, left=12pt, right=12pt,
        top=10pt, bottom=10pt]
      \renewcommand{\arraystretch}{1.35}%
      \begin{tabular}{@{}r@{\hspace{1.2em}}p{0.62\textwidth}@{}}
        #3
      \end{tabular}
    \end{tcolorbox}
    \vfill
    {\large #4\par}
    \vspace{0.3cm}
    {#5\par}
  \end{titlepage}
  \tableofcontents
  \newpage
}

%      \begin{document} ... \end{document}
%
%  Moteur : pdflatex (compiler deux fois pour la table des matières, ou
%  utiliser latexmk -pdf).
% =====================================================================

\usepackage[T1]{fontenc}
\usepackage[utf8]{inputenc}
\usepackage{lmodern}
\usepackage[french]{babel}
\usepackage{csquotes}
\frenchsetup{StandardItemLabels=true}
\usepackage{amsmath}

\usepackage[margin=2.5cm]{geometry}
\usepackage{graphicx}
\usepackage{float}
\usepackage{xcolor}
\usepackage{array}
\usepackage{booktabs}
\usepackage{enumitem}
\usepackage{listings}
\usepackage{tcolorbox}
\tcbuselibrary{skins,breakable}
\usepackage[hidelinks]{hyperref}

% ---------------------------------------------------------------------
%  Palette
% ---------------------------------------------------------------------
\definecolor{codebg}{RGB}{245,246,248}
\definecolor{codeframe}{RGB}{210,214,220}
\definecolor{codekw}{RGB}{170,40,60}
\definecolor{codestr}{RGB}{30,110,70}
\definecolor{codecmt}{RGB}{120,125,135}
\definecolor{accent}{RGB}{30,70,120}
\definecolor{warnbg}{RGB}{255,247,230}
\definecolor{warnframe}{RGB}{224,168,64}
\definecolor{keybg}{RGB}{235,240,248}
\definecolor{keyframe}{RGB}{120,150,190}

% ---------------------------------------------------------------------
%  Listings (blocs de code encadrés, monospace)
% ---------------------------------------------------------------------
\lstdefinestyle{report}{
  backgroundcolor=\color{codebg},
  frame=single,
  rulecolor=\color{codeframe},
  basicstyle=\ttfamily\small,
  keywordstyle=\color{codekw}\bfseries,
  stringstyle=\color{codestr},
  commentstyle=\color{codecmt}\itshape,
  numbers=none,
  breaklines=true,
  breakatwhitespace=false,
  showstringspaces=false,
  columns=fullflexible,
  keepspaces=true,
  xleftmargin=8pt,
  xrightmargin=8pt,
  aboveskip=10pt,
  belowskip=10pt,
  literate=%
    {é}{{\'e}}1 {è}{{\`e}}1 {ê}{{\^e}}1 {à}{{\`a}}1 {ç}{{\c c}}1
    {É}{{\'E}}1 {î}{{\^i}}1 {ï}{{\"i}}1 {ô}{{\^o}}1 {û}{{\^u}}1
    {ù}{{\`u}}1 {œ}{{\oe}}1 {»}{{\guillemotright}}1 {«}{{\guillemotleft}}1
}
\lstset{style=report}

% ---------------------------------------------------------------------
%  Captures d'écran : image + légende en italique (sans « Figure N »)
%  \shot{fichier.png}{légende}
%  Si l'image manque, un cadre de remplacement est affiché.
% ---------------------------------------------------------------------
\newcommand{\shot}[2]{%
  \begin{center}%
    \IfFileExists{images/#1}%
      {\includegraphics[width=0.92\linewidth]{images/#1}}%
      {\setlength{\fboxsep}{0pt}\fcolorbox{codeframe}{codebg}{%
         \parbox[c][3.2cm][c]{0.92\linewidth}{\centering\ttfamily\small
         images/#1\\[4pt]\normalfont\itshape(capture à ajouter)}}}%
    \\[5pt]{\itshape\small #2}%
  \end{center}%
}

% ---------------------------------------------------------------------
%  Encadrés « points clés » (bleu) et « attention » (orange)
% ---------------------------------------------------------------------
\newtcolorbox{keybox}[1][]{
  breakable, colback=keybg, colframe=keyframe, boxrule=0.6pt,
  arc=2pt, left=8pt, right=8pt, top=6pt, bottom=6pt, #1}
\newtcolorbox{warnbox}[1][]{
  breakable, colback=warnbg, colframe=warnframe, boxrule=0.6pt,
  arc=2pt, left=8pt, right=8pt, top=6pt, bottom=6pt, #1}

\setlist[itemize]{leftmargin=1.4em, itemsep=2pt, topsep=4pt}

% ---------------------------------------------------------------------
%  Page de garde + table des matières
%  \makecover{titre}{sous-titre}{lignes méta}{auteur}{date}
%  Les lignes méta s'écrivent avec \metarow{Étiquette}{Valeur}.
% ---------------------------------------------------------------------
\newcommand{\metarow}[2]{\textbf{#1} & #2 \\}
\newcommand{\makecover}[5]{%
  \begin{titlepage}
    \centering
    \vspace*{1.2cm}
    {\large\scshape Rapport de sécurité}\par
    \vspace{1.6cm}
    {\huge\bfseries #1\par}
    \vspace{0.5cm}
    {\Large #2\par}
    \vspace{2.2cm}
    \begin{tcolorbox}[width=0.86\textwidth, colback=codebg,
        colframe=codeframe, boxrule=0.8pt, arc=3pt, left=12pt, right=12pt,
        top=10pt, bottom=10pt]
      \renewcommand{\arraystretch}{1.35}%
      \begin{tabular}{@{}r@{\hspace{1.2em}}p{0.62\textwidth}@{}}
        #3
      \end{tabular}
    \end{tcolorbox}
    \vfill
    {\large #4\par}
    \vspace{0.3cm}
    {#5\par}
  \end{titlepage}
  \tableofcontents
  \newpage
}

%      \begin{document} ... \end{document}
%
%  Moteur : pdflatex (compiler deux fois pour la table des matières, ou
%  utiliser latexmk -pdf).
% =====================================================================

\usepackage[T1]{fontenc}
\usepackage[utf8]{inputenc}
\usepackage{lmodern}
\usepackage[french]{babel}
\usepackage{csquotes}
\frenchsetup{StandardItemLabels=true}
\usepackage{amsmath}

\usepackage[margin=2.5cm]{geometry}
\usepackage{graphicx}
\usepackage{float}
\usepackage{xcolor}
\usepackage{array}
\usepackage{booktabs}
\usepackage{enumitem}
\usepackage{listings}
\usepackage{tcolorbox}
\tcbuselibrary{skins,breakable}
\usepackage[hidelinks]{hyperref}

% ---------------------------------------------------------------------
%  Palette
% ---------------------------------------------------------------------
\definecolor{codebg}{RGB}{245,246,248}
\definecolor{codeframe}{RGB}{210,214,220}
\definecolor{codekw}{RGB}{170,40,60}
\definecolor{codestr}{RGB}{30,110,70}
\definecolor{codecmt}{RGB}{120,125,135}
\definecolor{accent}{RGB}{30,70,120}
\definecolor{warnbg}{RGB}{255,247,230}
\definecolor{warnframe}{RGB}{224,168,64}
\definecolor{keybg}{RGB}{235,240,248}
\definecolor{keyframe}{RGB}{120,150,190}

% ---------------------------------------------------------------------
%  Listings (blocs de code encadrés, monospace)
% ---------------------------------------------------------------------
\lstdefinestyle{report}{
  backgroundcolor=\color{codebg},
  frame=single,
  rulecolor=\color{codeframe},
  basicstyle=\ttfamily\small,
  keywordstyle=\color{codekw}\bfseries,
  stringstyle=\color{codestr},
  commentstyle=\color{codecmt}\itshape,
  numbers=none,
  breaklines=true,
  breakatwhitespace=false,
  showstringspaces=false,
  columns=fullflexible,
  keepspaces=true,
  xleftmargin=8pt,
  xrightmargin=8pt,
  aboveskip=10pt,
  belowskip=10pt,
  literate=%
    {é}{{\'e}}1 {è}{{\`e}}1 {ê}{{\^e}}1 {à}{{\`a}}1 {ç}{{\c c}}1
    {É}{{\'E}}1 {î}{{\^i}}1 {ï}{{\"i}}1 {ô}{{\^o}}1 {û}{{\^u}}1
    {ù}{{\`u}}1 {œ}{{\oe}}1 {»}{{\guillemotright}}1 {«}{{\guillemotleft}}1
}
\lstset{style=report}

% ---------------------------------------------------------------------
%  Captures d'écran : image + légende en italique (sans « Figure N »)
%  \shot{fichier.png}{légende}
%  Si l'image manque, un cadre de remplacement est affiché.
% ---------------------------------------------------------------------
\newcommand{\shot}[2]{%
  \begin{center}%
    \IfFileExists{images/#1}%
      {\includegraphics[width=0.92\linewidth]{images/#1}}%
      {\setlength{\fboxsep}{0pt}\fcolorbox{codeframe}{codebg}{%
         \parbox[c][3.2cm][c]{0.92\linewidth}{\centering\ttfamily\small
         images/#1\\[4pt]\normalfont\itshape(capture à ajouter)}}}%
    \\[5pt]{\itshape\small #2}%
  \end{center}%
}

% ---------------------------------------------------------------------
%  Encadrés « points clés » (bleu) et « attention » (orange)
% ---------------------------------------------------------------------
\newtcolorbox{keybox}[1][]{
  breakable, colback=keybg, colframe=keyframe, boxrule=0.6pt,
  arc=2pt, left=8pt, right=8pt, top=6pt, bottom=6pt, #1}
\newtcolorbox{warnbox}[1][]{
  breakable, colback=warnbg, colframe=warnframe, boxrule=0.6pt,
  arc=2pt, left=8pt, right=8pt, top=6pt, bottom=6pt, #1}

\setlist[itemize]{leftmargin=1.4em, itemsep=2pt, topsep=4pt}

% ---------------------------------------------------------------------
%  Page de garde + table des matières
%  \makecover{titre}{sous-titre}{lignes méta}{auteur}{date}
%  Les lignes méta s'écrivent avec \metarow{Étiquette}{Valeur}.
% ---------------------------------------------------------------------
\newcommand{\metarow}[2]{\textbf{#1} & #2 \\}
\newcommand{\makecover}[5]{%
  \begin{titlepage}
    \centering
    \vspace*{1.2cm}
    {\large\scshape Rapport de sécurité}\par
    \vspace{1.6cm}
    {\huge\bfseries #1\par}
    \vspace{0.5cm}
    {\Large #2\par}
    \vspace{2.2cm}
    \begin{tcolorbox}[width=0.86\textwidth, colback=codebg,
        colframe=codeframe, boxrule=0.8pt, arc=3pt, left=12pt, right=12pt,
        top=10pt, bottom=10pt]
      \renewcommand{\arraystretch}{1.35}%
      \begin{tabular}{@{}r@{\hspace{1.2em}}p{0.62\textwidth}@{}}
        #3
      \end{tabular}
    \end{tcolorbox}
    \vfill
    {\large #4\par}
    \vspace{0.3cm}
    {#5\par}
  \end{titlepage}
  \tableofcontents
  \newpage
}


\begin{document}

% --------------------------- Page de garde ---------------------------
\makecover
  {Titre du rapport}          % titre principal
  {sous-titre / précision}    % sous-titre
  {                            % lignes de l'encadré méta
    \metarow{Lab}{Nom du lab ou du challenge}
    \metarow{Classe}{CWE-XXX / OWASP AYY (2021), Catégorie}
    \metarow{Difficulté}{Practitioner / Expert / …}
    \metarow{Cible}{Composant visé}
    \metarow{Plateforme}{PortSwigger / Root-Me / …}
  }
  {Lucas Fagioli}             % auteur
  {1er janvier 2026}          % date

% ------------------------------ Corps --------------------------------
\section{Contexte}
Présentation de l'application, de la faille et de l'objectif.

\begin{itemize}
  \item \textbf{Point d'injection} : \dots
  \item \textbf{Objectif} : \dots
\end{itemize}

\subsection{Environnement}
\renewcommand{\arraystretch}{1.3}
\begin{tabular}{@{}p{0.30\textwidth}p{0.62\textwidth}@{}}
  \textbf{Système}        & \dots \\
  \textbf{Outil principal}& \dots \\
  \textbf{Cible}          & \dots \\
\end{tabular}

\shot{01-exemple.png}{Légende de la capture, en italique.}

\section{Reconnaissance}
\dots

\section{Exploitation}
\dots

% Bloc de code encadré :
\begin{lstlisting}
exemple de charge utile
\end{lstlisting}

% Encadré « point clé » :
\begin{keybox}
\textbf{Idée clé.} \dots
\end{keybox}

% Encadré « attention » :
\begin{warnbox}
\textbf{Attention.} \dots
\end{warnbox}

\section{Impact}
\dots\ Gravité indicative : \textbf{élevée}.

\section{Remédiation}
\begin{itemize}
  \item \textbf{Correctif principal} : \dots
\end{itemize}

\vspace{4pt}
\noindent\emph{Références} : \dots\ ; CWE-XXX.

\end{document}
